% ==============================================================================
% PAGE 37: CONSTRUCTION - PART 2 (Numbered Page 29)
% ==============================================================================
\noindent
\begin{spacing}{1.25}
{\fontsize{10.2}{14.5}\selectfont
The parabolic dish is then mounted on the top of the PVC shaft using a pulley and stud arrangement. The parabolic dish is made of satellite dish which is light in weight , and its surface is coated with reflective mirror pices to improve the reflection of sunlight. The shape of the dish is maintained accurately so that all reflected rays meet at one focal point. A copper tube painted in blackbody colour is fixed at this focal point using a supporting frame. This tube is connected with inlet and outlet pipes for water flow. To ensure smooth and stable rotation of the dish, a balancing weight is attached to the stud nut on the opposite side of the dish. This dead weight helps maintain proper balance of the parabolic reflector and prevents vibration or tilting during movement. The entire receiver assembly is held in place with metal nut \& bolts , and insulation material is added back side the tube to prevent heat loss.\\[4mm]
Next, the solar panel and battery system are installed on the same base frame. A 12V solar panel is fixed at a suitable angle to receive maximum sunlight. The panel output is connected to a 6V rechargeable battery through a charge controller. A PMDC motor is mounted on the base near to the shaft and connected to it using a V-belt and pulley arrangement of 25 cm diameter. Two LDR sensors are fixed on both sides of the parabolic dish, and all wiring between the panel, battery, sensors, and motor is neatly arranged using insulation tape and clips. After completing the assembly, the entire setup is checked for proper alignment, balance, and secure fitting of each component, ensuring the system is mechanically stable and ready for operation.
}
\end{spacing}
\clearpage
